\documentclass{beamer}
\usetheme{Madrid}
\usecolortheme{default}
\usepackage[utf8]{inputenc}
\usepackage{amsmath}
\usepackage{amsfonts}
\usepackage{amssymb}
\usepackage{tikz}
\usepackage{xcolor}
\usepackage{hyperref}

\title{CRISC Exam - Domain 4: Technology \& Security}
\subtitle{Technology Risk Questions with Explanations}
\author{Compiled from ISACA CRISC Review Manual 8th Edition}
\date{\today}

\begin{document}
	
	\begin{frame}
		\titlepage
	\end{frame}
	
	\begin{frame}{Outline}
		\tableofcontents
	\end{frame}
	
	% ================ SECTION 1: TECHNOLOGY ROADMAPS & ENTERPRISE ARCHITECTURE ================
	
	\section{Technology Roadmaps \& Enterprise Architecture}
	
	\begin{frame}{Q1: Purpose of Enterprise Architecture}
		\begin{block}{Question}
			What is the PRIMARY purpose of developing an enterprise architecture (EA)?
		\end{block}
		\begin{enumerate}[(A)]
			\item To document all hardware and software assets
			\item To deliver a view of the current state of IT, establish a future-state vision, and generate a strategy to move between them
			\item To replace the need for a risk assessment
			\item To satisfy regulatory audit requirements
		\end{enumerate}
	\end{frame}
	
	\begin{frame}{Q1: Explanation}
		\begin{block}{Correct Answer: B}
			Enterprise architecture delivers a view of the current state of IT, establishes a vision for a future state, and generates a strategy to move from current to future conditions that minimizes business disruption.
		\end{block}
		\begin{block}{Option Analysis}
			\begin{itemize}
				\item \textbf{A - Incorrect:} An asset inventory is only one component
				\item \textbf{B - Correct:} EA frames how information enables the organization
				\item \textbf{C - Incorrect:} EA complements, not replaces, risk assessment
				\item \textbf{D - Incorrect:} Compliance is a secondary benefit
			\end{itemize}
		\end{block}
	\end{frame}
	
	\begin{frame}{Q2: EA Framework Key Topics}
		\begin{block}{Question}
			All EA frameworks provide structured guidance across which four key topics?
		\end{block}
		\begin{enumerate}[(A)]
			\item Budgeting, staffing, training, and auditing
			\item Documentation, notation, process, and organization
			\item Hardware, software, networks, and data
			\item Identify, assess, respond, and monitor
		\end{enumerate}
	\end{frame}
	
	\begin{frame}{Q2: Explanation}
		\begin{block}{Correct Answer: B}
			All EA frameworks provide structured guidance across four key topics: documentation, notation, process, and organization.
		\end{block}
		\begin{block}{Option Analysis}
			\begin{itemize}
				\item \textbf{A - Incorrect:} These are management activities
				\item \textbf{B - Correct:} Documentation, notation, process, organization
				\item \textbf{C - Incorrect:} These are technology domains
				\item \textbf{D - Incorrect:} These are risk management phases
			\end{itemize}
		\end{block}
	\end{frame}
	
	\begin{frame}{Q3: Technology Roadmap Purpose}
		\begin{block}{Question}
			What is the purpose of a technology roadmap?
		\end{block}
		\begin{enumerate}[(A)]
			\item To list all current IT projects
			\item To link short- and long-term business goals with specific technology solutions
			\item To document the IT organizational chart
			\item To record all IT incidents
		\end{enumerate}
	\end{frame}
	
	\begin{frame}{Q3: Explanation}
		\begin{block}{Correct Answer: B}
			A technology roadmap is the planning process that links new short- and long-term goals identified by the EA with specific technology solutions.
		\end{block}
		\begin{block}{Option Analysis}
			\begin{itemize}
				\item \textbf{A - Incorrect:} A project list is not a roadmap
				\item \textbf{B - Correct:} Roadmaps link goals to technology solutions
				\item \textbf{C - Incorrect:} Org charts are unrelated to roadmaps
				\item \textbf{D - Incorrect:} Incident records are operational data
			\end{itemize}
		\end{block}
	\end{frame}
	
	\begin{frame}{Q4: Risk Without Enterprise Architecture}
		\begin{block}{Question}
			In environments where an EA is either immature or absent, what should the risk practitioner do?
		\end{block}
		\begin{enumerate}[(A)]
			\item Suspend all risk management activities until an EA is built
			\item Place greater emphasis on technology-specific assessments to build a piecemeal vision of current IT risk
			\item Rely solely on qualitative risk assessments
			\item Delegate all risk assessment to the IT department
		\end{enumerate}
	\end{frame}
	
	\begin{frame}{Q4: Explanation}
		\begin{block}{Correct Answer: B}
			Where an EA is immature or absent, the risk practitioner must place greater emphasis on technology-specific assessments to build a piecemeal vision of current IT risk.
		\end{block}
		\begin{block}{Option Analysis}
			\begin{itemize}
				\item \textbf{A - Incorrect:} Risk management cannot be suspended
				\item \textbf{B - Correct:} Technology-specific assessments fill the gap
				\item \textbf{C - Incorrect:} Qualitative assessment alone is insufficient
				\item \textbf{D - Incorrect:} Risk management is an enterprise responsibility
			\end{itemize}
		\end{block}
	\end{frame}
	
	\begin{frame}{Q5: Maturity Model Purpose}
		\begin{block}{Question}
			What is the goal of a maturity model in the context of enterprise architecture?
		\end{block}
		\begin{enumerate}[(A)]
			\item To dictate exactly how processes should be performed
			\item To define what outcomes are desirable and help organizations gauge progress against published baselines
			\item To replace internal audits
			\item To certify compliance with regulations
		\end{enumerate}
	\end{frame}
	
	\begin{frame}{Q5: Explanation}
		\begin{block}{Correct Answer: B}
			The goal of maturity models is not to dictate how something should be done but rather what outcomes are desirable, helping organizations gauge their progress.
		\end{block}
		\begin{block}{Option Analysis}
			\begin{itemize}
				\item \textbf{A - Incorrect:} Maturity models focus on outcomes, not prescriptive steps
				\item \textbf{B - Correct:} They define desirable outcomes
				\item \textbf{C - Incorrect:} They do not replace audits
				\item \textbf{D - Incorrect:} They are not certification mechanisms
			\end{itemize}
		\end{block}
	\end{frame}
	
	% ================ SECTION 2: OPERATIONS MANAGEMENT ================
	
	\section{Operations Management: Hardware, Software \& Networks}
	
	\begin{frame}{Q6: Hardware Risk Areas}
		\begin{block}{Question}
			Which of the following is an area of concern associated with hardware?
		\end{block}
		\begin{enumerate}[(A)]
			\item Logic flaws or semantic errors
			\item Use of default credentials and unnecessary services
			\item Inadequate data validation
			\item Lack of version control
		\end{enumerate}
	\end{frame}
	
	\begin{frame}{Q6: Explanation}
		\begin{block}{Correct Answer: B}
			Hardware concerns include obsolescence, poor maintenance, misconfiguration, use of default credentials, unnecessary services, physical loss, insecure disposal, and component failure.
		\end{block}
		\begin{block}{Option Analysis}
			\begin{itemize}
				\item \textbf{A - Incorrect:} Logic flaws are a software concern
				\item \textbf{B - Correct:} Default credentials are a hardware concern
				\item \textbf{C - Incorrect:} Data validation is a software concern
				\item \textbf{D - Incorrect:} Version control is a software concern
			\end{itemize}
		\end{block}
	\end{frame}
	
	\begin{frame}{Q7: Supply Chain Hardware Risk}
		\begin{block}{Question}
			What should organizations do to mitigate the risk of purchasing hardware infected with security vulnerabilities during manufacturing or delivery?
		\end{block}
		\begin{enumerate}[(A)]
			\item Purchase only from the lowest-cost vendor
			\item Use trusted vendors or suppliers and validate maintenance hooks
			\item Rely solely on vendor self-attestation
			\item Avoid all hardware purchases
		\end{enumerate}
	\end{frame}
	
	\begin{frame}{Q7: Explanation}
		\begin{block}{Correct Answer: B}
			Risk practitioners should encourage use of trusted vendors and validate that maintenance hooks are documented and either secured or eliminated.
		\end{block}
		\begin{block}{Option Analysis}
			\begin{itemize}
				\item \textbf{A - Incorrect:} Lowest cost may increase risk
				\item \textbf{B - Correct:} Trusted vendors and validated hooks
				\item \textbf{C - Incorrect:} Self-attestation is insufficient
				\item \textbf{D - Incorrect:} Avoiding all hardware is impractical
			\end{itemize}
		\end{block}
	\end{frame}
	
	\begin{frame}{Q8: Operating System Risk Areas}
		\begin{block}{Question}
			Which of the following is a concern specific to operating systems?
		\end{block}
		\begin{enumerate}[(A)]
			\item Kernel-level vulnerabilities such as rootkits and hypervisor escapes
			\item Poor physical access controls
			\item Inadequate HVAC systems
			\item Cable exceeding approved length
		\end{enumerate}
	\end{frame}
	
	\begin{frame}{Q8: Explanation}
		\begin{block}{Correct Answer: A}
			OS concerns include unpatched vulnerabilities, buffer overflows, misconfiguration, weak access controls, and kernel-level vulnerabilities such as rootkits and hypervisor escapes.
		\end{block}
		\begin{block}{Option Analysis}
			\begin{itemize}
				\item \textbf{A - Correct:} Kernel-level vulnerabilities are OS concerns
				\item \textbf{B - Incorrect:} Physical access is an environmental control
				\item \textbf{C - Incorrect:} HVAC is an environmental control
				\item \textbf{D - Incorrect:} Cabling is a network concern
			\end{itemize}
		\end{block}
	\end{frame}
	
	\begin{frame}{Q9: Patch Management Lesson}
		\begin{block}{Question}
			The 2024 CrowdStrike global outage best illustrates which risk?
		\end{block}
		\begin{enumerate}[(A)]
			\item The need to avoid all software updates
			\item The importance of not skipping testing and patch management processes
			\item The risk of using cloud services
			\item The danger of open-source software
		\end{enumerate}
	\end{frame}
	
	\begin{frame}{Q9: Explanation}
		\begin{block}{Correct Answer: B}
			Root cause analysis discovered the update had not undergone routine patch management processes. Testing should not be skipped to expedite patching.
		\end{block}
		\begin{block}{Option Analysis}
			\begin{itemize}
				\item \textbf{A - Incorrect:} Avoiding updates leaves systems vulnerable
				\item \textbf{B - Correct:} Testing before deployment is essential
				\item \textbf{C - Incorrect:} This was not cloud-specific
				\item \textbf{D - Incorrect:} This was a commercial product
			\end{itemize}
		\end{block}
	\end{frame}
	
	\begin{frame}{Q10: Application Risk from Updates}
		\begin{block}{Question}
			In mobile computing, why can application updates introduce risk?
		\end{block}
		\begin{enumerate}[(A)]
			\item Updates always improve security
			\item A major version release may change or remove functionality, and vendors may end support for prior versions
			\item Updates are always optional
			\item Updates never affect functionality
		\end{enumerate}
	\end{frame}
	
	\begin{frame}{Q10: Explanation}
		\begin{block}{Correct Answer: B}
			A major version release may change or remove functionality. Vendors may end support for prior versions, causing sudden loss of features that impact production.
		\end{block}
		\begin{block}{Option Analysis}
			\begin{itemize}
				\item \textbf{A - Incorrect:} Updates can introduce new risk
				\item \textbf{B - Correct:} Functionality may change or be removed
				\item \textbf{C - Incorrect:} Updates may be effectively mandatory
				\item \textbf{D - Incorrect:} Updates frequently affect functionality
			\end{itemize}
		\end{block}
	\end{frame}
	
	\begin{frame}{Q11: Database Security Controls}
		\begin{block}{Question}
			Which of the following is a method for providing database security?
		\end{block}
		\begin{enumerate}[(A)]
			\item Placing database servers directly on the internet
			\item Encryption of sensitive data in the database and restricting administrator-level access
			\item Disabling all backups
			\item Allowing unrestricted access to all tables
		\end{enumerate}
	\end{frame}
	
	\begin{frame}{Q11: Explanation}
		\begin{block}{Correct Answer: B}
			Database security is provided through encryption of sensitive data, use of database views, secure protocols, content-based access controls, restricting administrator-level access, and backups. Database servers should not be directly accessible from outside.
		\end{block}
		\begin{block}{Option Analysis}
			\begin{itemize}
				\item \textbf{A - Incorrect:} Direct internet exposure is a major risk
				\item \textbf{B - Correct:} Encryption and restricted admin access
				\item \textbf{C - Incorrect:} Backups are essential
				\item \textbf{D - Incorrect:} Unrestricted access violates least privilege
			\end{itemize}
		\end{block}
	\end{frame}
	
	\begin{frame}{Q12: Software Utilities Risk}
		\begin{block}{Question}
			Why are compromised device drivers a significant risk?
		\end{block}
		\begin{enumerate}[(A)]
			\item They only affect printer functionality
			\item They can alter intended commands to yield different results or fail to deliver instructions
			\item They are always digitally signed and trusted
			\item They have no interaction with hardware
		\end{enumerate}
	\end{frame}
	
	\begin{frame}{Q12: Explanation}
		\begin{block}{Correct Answer: B}
			Compromised drivers may interfere with proper operations of hardware by altering intended commands to yield different results or failing to deliver instructions.
		\end{block}
		\begin{block}{Option Analysis}
			\begin{itemize}
				\item \textbf{A - Incorrect:} Drivers affect all hardware
				\item \textbf{B - Correct:} Compromised drivers can alter commands
				\item \textbf{C - Incorrect:} Digital signing does not guarantee safety
				\item \textbf{D - Incorrect:} Drivers bridge software and hardware
			\end{itemize}
		\end{block}
	\end{frame}
	
	\begin{frame}{Q13: Environmental Controls - Power}
		\begin{block}{Question}
			Which of the following is an area of concern associated with power environmental controls?
		\end{block}
		\begin{enumerate}[(A)]
			\item Clogged filters
			\item Surges, spikes, sags, brownouts, and generator latency
			\item High humidity causing corrosion
			\item Low humidity causing static
		\end{enumerate}
	\end{frame}
	
	\begin{frame}{Q13: Explanation}
		\begin{block}{Correct Answer: B}
			Power concerns include surges, spikes, sags, brownouts, and faults, plus generator latency or spin-up time.
		\end{block}
		\begin{block}{Option Analysis}
			\begin{itemize}
				\item \textbf{A - Incorrect:} Clogged filters are an HVAC concern
				\item \textbf{B - Correct:} Power effects
				\item \textbf{C - Incorrect:} High humidity is an HVAC concern
				\item \textbf{D - Incorrect:} Low humidity is an HVAC concern
			\end{itemize}
		\end{block}
	\end{frame}
	
	\begin{frame}{Q14: Firewall Configuration Review}
		\begin{block}{Question}
			Which provides the BEST assurance that a firewall is configured in compliance with an enterprise's security policy?
		\end{block}
		\begin{enumerate}[(A)]
			\item Reviewing actual procedures
			\item Interviewing the firewall administrator
			\item Reviewing parameter settings
			\item Reviewing the log file for recent attacks
		\end{enumerate}
	\end{frame}
	
	\begin{frame}{Q14: Explanation}
		\begin{block}{Correct Answer: C}
			A review of the parameter settings provides a good basis for comparison of the actual configuration to the security policy.
		\end{block}
		\begin{block}{Option Analysis}
			\begin{itemize}
				\item \textbf{A - Incorrect:} Procedures describe intent
				\item \textbf{B - Incorrect:} Interviews do not confirm configuration
				\item \textbf{C - Correct:} Parameter settings show actual configuration
				\item \textbf{D - Incorrect:} Logs show activity, not configuration
			\end{itemize}
		\end{block}
	\end{frame}
	
	\begin{frame}{Q15: Firewall Log Review}
		\begin{block}{Question}
			Why should firewall logs be regularly reviewed?
		\end{block}
		\begin{enumerate}[(A)]
			\item To detect suspicious activity and verify logging is functioning
			\item To update the firewall's configuration automatically
			\item To replace the need for an intrusion detection system
			\item To calculate the firewall's annual loss expectancy
		\end{enumerate}
	\end{frame}
	
	\begin{frame}{Q15: Explanation}
		\begin{block}{Correct Answer: A}
			Firewall logs should be regularly reviewed to detect any suspicious activity.
		\end{block}
		\begin{block}{Option Analysis}
			\begin{itemize}
				\item \textbf{A - Correct:} Logs are reviewed to detect suspicious activity
				\item \textbf{B - Incorrect:} Log review does not update configuration
				\item \textbf{C - Incorrect:} Logs complement, not replace, IDS
				\item \textbf{D - Incorrect:} ALE is calculated from asset value
			\end{itemize}
		\end{block}
	\end{frame}
	
	\begin{frame}{Q16: Proxy Function}
		\begin{block}{Question}
			What is the role of a proxy in a network?
		\end{block}
		\begin{enumerate}[(A)]
			\item To act as an intermediary between two communicating parties
			\item To permanently store all network traffic
			\item To replace the need for a firewall
			\item To assign IP addresses to network devices
		\end{enumerate}
	\end{frame}
	
	\begin{frame}{Q16: Explanation}
		\begin{block}{Correct Answer: A}
			A proxy acts as an intermediary, presenting itself to each side as if it were the other end. This allows filtering and examination of suspicious activity.
		\end{block}
		\begin{block}{Option Analysis}
			\begin{itemize}
				\item \textbf{A - Correct:} Intermediary
				\item \textbf{B - Incorrect:} Proxies filter, not store
				\item \textbf{C - Incorrect:} Proxies complement firewalls
				\item \textbf{D - Incorrect:} DHCP assigns IP addresses
			\end{itemize}
		\end{block}
	\end{frame}
	
	\begin{frame}{Q17: IDS vs. IPS}
		\begin{block}{Question}
			What is the key difference between an IDS and an IPS?
		\end{block}
		\begin{enumerate}[(A)]
			\item An IDS is more expensive than an IPS
			\item An IDS detects and alerts, while an IPS detects and blocks traffic in-line
			\item An IPS has no effect on network throughput
			\item An IDS must be placed in the line of traffic
		\end{enumerate}
	\end{frame}
	
	\begin{frame}{Q17: Explanation}
		\begin{block}{Correct Answer: B}
			An IDS alerts without affecting throughput. An IPS combines detection with blocking and must be placed in-line, generally slowing the network.
		\end{block}
		\begin{block}{Option Analysis}
			\begin{itemize}
				\item \textbf{A - Incorrect:} Cost is not the differentiator
				\item \textbf{B - Correct:} IDS alerts; IPS blocks in-line
				\item \textbf{C - Incorrect:} An IPS generally slows throughput
				\item \textbf{D - Incorrect:} An IPS, not an IDS, is in-line
			\end{itemize}
		\end{block}
	\end{frame}
	
	\begin{frame}{Q18: DNS Risk}
		\begin{block}{Question}
			Which of the following is a security risk associated with the Domain Name System (DNS)?
		\end{block}
		\begin{enumerate}[(A)]
			\item DNS replies can be used in amplification attacks to flood a victim's system
			\item DNS encrypts all network traffic
			\item DNS replaces the need for IP addressing
			\item DNS prevents all man-in-the-middle attacks
		\end{enumerate}
	\end{frame}
	
	\begin{frame}{Q18: Explanation}
		\begin{block}{Correct Answer: A}
			False DNS replies can misroute traffic, and DNS replies have been used in amplification attacks. Attacks against DNS servers can impede name resolution.
		\end{block}
		\begin{block}{Option Analysis}
			\begin{itemize}
				\item \textbf{A - Correct:} DNS amplification attacks are a known risk
				\item \textbf{B - Incorrect:} DNS does not encrypt traffic
				\item \textbf{C - Incorrect:} DNS maps names to IP addresses
				\item \textbf{D - Incorrect:} DNS does not prevent MITM attacks
			\end{itemize}
		\end{block}
	\end{frame}
	
	\begin{frame}{Q19: Rogue Wireless Access Points}
		\begin{block}{Question}
			Why are rogue wireless access points a significant vulnerability?
		\end{block}
		\begin{enumerate}[(A)]
			\item They improve network performance
			\item They may permit a person to access the network directly without going through a secure device or bypassing firewalls
			\item They are always detected by antivirus software
			\item They require no physical access to install
		\end{enumerate}
	\end{frame}
	
	\begin{frame}{Q19: Explanation}
		\begin{block}{Correct Answer: B}
			Rogue access points may permit a person to access the network directly, bypassing firewalls or other layers of defense. Wireless discovery scans should be part of the vulnerability assessment process.
		\end{block}
		\begin{block}{Option Analysis}
			\begin{itemize}
				\item \textbf{A - Incorrect:} They do not improve security
				\item \textbf{B - Correct:} They bypass network security controls
				\item \textbf{C - Incorrect:} Antivirus does not detect rogue APs
				\item \textbf{D - Incorrect:} They usually require physical access
			\end{itemize}
		\end{block}
	\end{frame}
	
	\begin{frame}{Q20: Mesh Topology}
		\begin{block}{Question}
			Which network topology provides the highest availability because traffic can route around a failure?
		\end{block}
		\begin{enumerate}[(A)]
			\item Bus
			\item Star
			\item Ring
			\item Mesh
		\end{enumerate}
	\end{frame}
	
	\begin{frame}{Q20: Explanation}
		\begin{block}{Correct Answer: D}
			In a mesh topology, many devices are connected to many other devices so traffic can route around a failure. The internet is a partial mesh.
		\end{block}
		\begin{block}{Option Analysis}
			\begin{itemize}
				\item \textbf{A - Incorrect:} A cut cable in a bus fails all downstream users
				\item \textbf{B - Incorrect:} Loss of central switch in a star affects all
				\item \textbf{C - Incorrect:} A single cut affects a unidirectional ring
				\item \textbf{D - Correct:} Mesh routes around failures
			\end{itemize}
		\end{block}
	\end{frame}
	
	\begin{frame}{Q21: LAN vs. WAN Protocols}
		\begin{block}{Question}
			Which protocol type is best suited for a WAN where traffic crosses network borders?
		\end{block}
		\begin{enumerate}[(A)]
			\item Interior gateway protocols (IGPs) such as OSPF
			\item Exterior gateway protocols (EGPs) such as BGP
			\item Routing Information Protocol (RIP)
			\item Enhanced Interior Gateway Routing Protocol (EIGRP)
		\end{enumerate}
	\end{frame}
	
	\begin{frame}{Q21: Explanation}
		\begin{block}{Correct Answer: B}
			WANs benefit from EGPs, which are more complex than IGPs. The most widely used EGP is BGP.
		\end{block}
		\begin{block}{Option Analysis}
			\begin{itemize}
				\item \textbf{A - Incorrect:} IGPs are used within a LAN
				\item \textbf{B - Correct:} EGPs such as BGP are used for WANs
				\item \textbf{C - Incorrect:} RIP is an IGP
				\item \textbf{D - Incorrect:} EIGRP is an IGP
			\end{itemize}
		\end{block}
	\end{frame}
	
	\begin{frame}{Q22: Software-Defined Networking}
		\begin{block}{Question}
			What is a key benefit of software-defined networking (SDN)?
		\end{block}
		\begin{enumerate}[(A)]
			\item It eliminates the need for network security controls
			\item It separates network control and forwarding functions, allowing dynamic shaping of traffic
			\item It permanently fixes network configurations
			\item It removes the need for a network administrator
		\end{enumerate}
	\end{frame}
	
	\begin{frame}{Q22: Explanation}
		\begin{block}{Correct Answer: B}
			SDN separates control and forwarding functions, allowing dynamic traffic shaping to meet business requirements on short notice.
		\end{block}
		\begin{block}{Option Analysis}
			\begin{itemize}
				\item \textbf{A - Incorrect:} SDN does not eliminate security needs
				\item \textbf{B - Correct:} Dynamic traffic shaping
				\item \textbf{C - Incorrect:} SDN enables dynamic configurations
				\item \textbf{D - Incorrect:} Administrators are still required
			\end{itemize}
		\end{block}
	\end{frame}
	
	\begin{frame}{Q23: Demilitarized Zone Purpose}
		\begin{block}{Question}
			What is the purpose of a demilitarized zone (DMZ) in a network?
		\end{block}
		\begin{enumerate}[(A)]
			\item To store critical organizational assets securely
			\item To house services intended for outside use, protecting internal systems
			\item To replace internal firewalls
			\item To host the organization's backup data
		\end{enumerate}
	\end{frame}
	
	\begin{frame}{Q23: Explanation}
		\begin{block}{Correct Answer: B}
			The DMZ houses externally facing services in a segment configured to allow open access. Critical assets are kept out of the DMZ.
		\end{block}
		\begin{block}{Option Analysis}
			\begin{itemize}
				\item \textbf{A - Incorrect:} Critical assets are kept out of the DMZ
				\item \textbf{B - Correct:} DMZ houses externally facing services
				\item \textbf{C - Incorrect:} DMZ complements firewalls
				\item \textbf{D - Incorrect:} Backup data should not be in the DMZ
			\end{itemize}
		\end{block}
	\end{frame}
	
	\begin{frame}{Q24: VPN Purpose}
		\begin{block}{Question}
			What is the primary purpose of a virtual private network (VPN)?
		\end{block}
		\begin{enumerate}[(A)]
			\item To make the network faster
			\item To logically separate traffic using encrypted tunnels, making it unreadable except by intended recipients
			\item To replace firewalls
			\item To assign IP addresses
		\end{enumerate}
	\end{frame}
	
	\begin{frame}{Q24: Explanation}
		\begin{block}{Correct Answer: B}
			A VPN logically separates data traffic by digitally constructed tunnels, configured with encryption so traffic is unreadable except by intended recipients.
		\end{block}
		\begin{block}{Option Analysis}
			\begin{itemize}
				\item \textbf{A - Incorrect:} VPNs add processing overhead
				\item \textbf{B - Correct:} VPNs create encrypted tunnels
				\item \textbf{C - Incorrect:} VPNs complement firewalls
				\item \textbf{D - Incorrect:} DHCP assigns IP addresses
			\end{itemize}
		\end{block}
	\end{frame}
	
	\begin{frame}{Q25: Zero Trust Architecture Principle}
		\begin{block}{Question}
			According to NIST SP 800-207, zero trust focuses on protecting:
		\end{block}
		\begin{enumerate}[(A)]
			\item Network segments
			\item Resources rather than network segments
			\item Only the perimeter firewall
			\item Only internal employees
		\end{enumerate}
	\end{frame}
	
	\begin{frame}{Q25: Explanation}
		\begin{block}{Correct Answer: B}
			Zero trust focuses on protecting resources (assets, services, workflows, accounts), not network segments, as the network location is no longer the prime component of the security posture.
		\end{block}
		\begin{block}{Option Analysis}
			\begin{itemize}
				\item \textbf{A - Incorrect:} Zero trust moves away from segment-based protection
				\item \textbf{B - Correct:} Focuses on protecting resources
				\item \textbf{C - Incorrect:} Perimeter-only protection is opposite of zero trust
				\item \textbf{D - Incorrect:} Applies to all users and assets
			\end{itemize}
		\end{block}
	\end{frame}
	
	\begin{frame}{Q26: Technology Refresh Considerations}
		\begin{block}{Question}
			Which of the following is a risk consideration related to technology refresh?
		\end{block}
		\begin{enumerate}[(A)]
			\item Age of equipment and availability of replacement parts
			\item The number of employees in the IT department
			\item The color scheme of the data center
			\item The brand of coffee served in the office
		\end{enumerate}
	\end{frame}
	
	\begin{frame}{Q26: Explanation}
		\begin{block}{Correct Answer: A}
			Risk considerations include age of equipment, available expertise, vendor viability, documentation, availability of replacement parts, ability to test, operating environment, patch ability, willingness to change, and vendor lock-in.
		\end{block}
		\begin{block}{Option Analysis}
			\begin{itemize}
				\item \textbf{A - Correct:} Age and parts availability
				\item \textbf{B - Incorrect:} Not a technology refresh factor
				\item \textbf{C - Incorrect:} Irrelevant
				\item \textbf{D - Incorrect:} Irrelevant
			\end{itemize}
		\end{block}
	\end{frame}
	
	\begin{frame}{Q27: Legacy System Risk}
		\begin{block}{Question}
			When legacy systems cannot be replaced, what is the risk practitioner's role?
		\end{block}
		\begin{enumerate}[(A)]
			\item Reject any course of action involving the legacy system
			\item Work with the risk owner to reach an acceptable level of risk using compensating controls
			\item Mandate immediate replacement regardless of cost
			\item Ignore the legacy system since it is unavoidable
		\end{enumerate}
	\end{frame}
	
	\begin{frame}{Q27: Explanation}
		\begin{block}{Correct Answer: B}
			Some legacy systems may not be replaceable. The risk practitioner's role is to work with the risk owner to reach an acceptable level of risk using compensating controls such as additional backups, segmentation, and increased monitoring.
		\end{block}
		\begin{block}{Option Analysis}
			\begin{itemize}
				\item \textbf{A - Incorrect:} Rejection is not the role
				\item \textbf{B - Correct:} Work with risk owner on compensating controls
				\item \textbf{C - Incorrect:} Mandating replacement ignores cost-benefit
				\item \textbf{D - Incorrect:} Ignoring increases risk
			\end{itemize}
		\end{block}
	\end{frame}
	
	\begin{frame}{Q28: IT Operations Security Problems}
		\begin{block}{Question}
			Which security problem is commonly seen within the IT department due to complacency?
		\end{block}
		\begin{enumerate}[(A)]
			\item Enforcing strict separation of duties
			\item Sharing high-level accounts and passwords and lack of control over changes to production systems
			\item Conducting regular vulnerability assessments
			\item Maintaining complete documentation
		\end{enumerate}
	\end{frame}
	
	\begin{frame}{Q28: Explanation}
		\begin{block}{Correct Answer: B}
			Where complacency reigns, risk ignorance can result in breaches. Common problems include sharing high-level accounts and passwords and lack of control over changes to production.
		\end{block}
		\begin{block}{Option Analysis}
			\begin{itemize}
				\item \textbf{A - Incorrect:} Good practice, not a problem
				\item \textbf{B - Correct:} Shared accounts and uncontrolled changes
				\item \textbf{C - Incorrect:} Regular assessments are good practice
				\item \textbf{D - Incorrect:} Documentation is good practice
			\end{itemize}
		\end{block}
	\end{frame}
	
	\begin{frame}{Q29: Configuration Management Goal}
		\begin{block}{Question}
			What is the goal of configuration management?
		\end{block}
		\begin{enumerate}[(A)]
			\item To have a single approved way in which systems or devices intended for a particular role should be configured
			\item To allow each administrator to configure systems as they prefer
			\item To eliminate all system changes
			\item To replace the need for change management
		\end{enumerate}
	\end{frame}
	
	\begin{frame}{Q29: Explanation}
		\begin{block}{Correct Answer: A}
			A single approved configuration makes it easier to deploy new systems, test patches and upgrades, identify malware, and manage the enterprise.
		\end{block}
		\begin{block}{Option Analysis}
			\begin{itemize}
				\item \textbf{A - Correct:} Single approved configuration
				\item \textbf{B - Incorrect:} Individual preference creates inconsistency
				\item \textbf{C - Incorrect:} Changes must be controlled, not eliminated
				\item \textbf{D - Incorrect:} Configuration management works with change management
			\end{itemize}
		\end{block}
	\end{frame}
	
	\begin{frame}{Q30: Hardening Process}
		\begin{block}{Question}
			What does the hardening process aim to achieve?
		\end{block}
		\begin{enumerate}[(A)]
			\item To install all available software
			\item To reduce vulnerability by limiting attack vectors, disabling unnecessary functionality, and changing default credentials
			\item To increase the number of open ports
			\item To enable all services by default
		\end{enumerate}
	\end{frame}
	
	\begin{frame}{Q30: Explanation}
		\begin{block}{Correct Answer: B}
			Hardening reduces vulnerability by limiting attack vectors. Unnecessary ports, services, and protocols are disabled; default accounts and passwords are changed or disabled.
		\end{block}
		\begin{block}{Option Analysis}
			\begin{itemize}
				\item \textbf{A - Incorrect:} Installing all software increases attack surface
				\item \textbf{B - Correct:} Hardening limits attack vectors
				\item \textbf{C - Incorrect:} Open ports increase vulnerability
				\item \textbf{D - Incorrect:} Unnecessary services should be disabled
			\end{itemize}
		\end{block}
	\end{frame}
	
	\begin{frame}{Q31: Cloud Computing Essential Characteristic}
		\begin{block}{Question}
			Which cloud characteristic allows capabilities to be rapidly and elastically provisioned to scale out and rapidly released to scale in?
		\end{block}
		\begin{enumerate}[(A)]
			\item On-demand self-service
			\item Broad network access
			\item Rapid elasticity
			\item Measured service
		\end{enumerate}
	\end{frame}
	
	\begin{frame}{Q31: Explanation}
		\begin{block}{Correct Answer: C}
			Rapid elasticity means capabilities can be rapidly and elastically provisioned and released. To the customer, capabilities often appear unlimited and can be purchased in any quantity at any time.
		\end{block}
		\begin{block}{Option Analysis}
			\begin{itemize}
				\item \textbf{A - Incorrect:} On-demand self-service is automatic provisioning
				\item \textbf{B - Incorrect:} Broad network access is accessibility
				\item \textbf{C - Correct:} Rapid elasticity
				\item \textbf{D - Incorrect:} Measured service is metering
			\end{itemize}
		\end{block}
	\end{frame}
	
	\begin{frame}{Q32: Cloud Migration Consideration}
		\begin{block}{Question}
			Which of the following is a common consideration before migrating to a cloud-based environment?
		\end{block}
		\begin{enumerate}[(A)]
			\item Increased visibility and control over operations
			\item Reduced visibility and control, unauthorized use and abuse, and provider's privileged insiders
			\item Guaranteed unlimited storage at no cost
			\item Elimination of all security responsibilities
		\end{enumerate}
	\end{frame}
	
	\begin{frame}{Q32: Explanation}
		\begin{block}{Correct Answer: B}
			Common cloud considerations include reduced visibility and control, unauthorized use and abuse, provider's privileged insiders, compromised APIs, multi-tenant separation failures, incomplete data deletion, stolen credentials, vendor lock-in, and insufficient due diligence.
		\end{block}
		\begin{block}{Option Analysis}
			\begin{itemize}
				\item \textbf{A - Incorrect:} Cloud generally reduces visibility
				\item \textbf{B - Correct:} Key cloud risk considerations
				\item \textbf{C - Incorrect:} Storage is not unlimited or free
				\item \textbf{D - Incorrect:} Organizations retain responsibilities
			\end{itemize}
		\end{block}
	\end{frame}
	
	\begin{frame}{Q33: Cloud Provider Certification Misconception}
		\begin{block}{Question}
			What common misconception do organizations have regarding cloud provider certifications?
		\end{block}
		\begin{enumerate}[(A)]
			\item Certifications are always too expensive
			\item A cloud provider's certification automatically transfers to the customer's environment
			\item Certifications guarantee 100\% uptime
			\item Certifications eliminate the need for customer controls
		\end{enumerate}
	\end{frame}
	
	\begin{frame}{Q33: Explanation}
		\begin{block}{Correct Answer: B}
			A common misconception is that if a cloud provider has a certification, that certification transfers to the customer's environment. This is not the case.
		\end{block}
		\begin{block}{Option Analysis}
			\begin{itemize}
				\item \textbf{A - Incorrect:} Cost is not the misconception
				\item \textbf{B - Correct:} Certifications do not transfer
				\item \textbf{C - Incorrect:} Certifications do not guarantee uptime
				\item \textbf{D - Incorrect:} Customer controls are still needed
			\end{itemize}
		\end{block}
	\end{frame}
	
	\begin{frame}{Q34: Big Data Risk - Amplified Impact}
		\begin{block}{Question}
			Which describes an amplified technical impact risk associated with big data?
		\end{block}
		\begin{enumerate}[(A)]
			\item Data is always more secure when centralized
			\item If an unauthorized user gains access to centralized repositories, the entirety of the data is in jeopardy
			\item Big data eliminates privacy concerns
			\item Big data reduces the need for data governance
		\end{enumerate}
	\end{frame}
	
	\begin{frame}{Q34: Explanation}
		\begin{block}{Correct Answer: B}
			Amplified technical impact occurs when an unauthorized user gains access to centralized repositories, putting the entirety of the data in jeopardy. Other risks include privacy concerns and governance challenges.
		\end{block}
		\begin{block}{Option Analysis}
			\begin{itemize}
				\item \textbf{A - Incorrect:} Centralization increases risk
				\item \textbf{B - Correct:} Centralized data magnifies breach impact
				\item \textbf{C - Incorrect:} Big data increases privacy concerns
				\item \textbf{D - Incorrect:} Big data increases governance need
			\end{itemize}
		\end{block}
	\end{frame}
	
	\begin{frame}{Q35: Large Language Model Risk}
		\begin{block}{Question}
			What is a key risk associated with large language models (LLMs) and generative AI?
		\end{block}
		\begin{enumerate}[(A)]
			\item They always provide accurate information
			\item They can create false information or "hallucinations"
			\item They eliminate the need for human oversight
			\item They cannot be used for malicious purposes
		\end{enumerate}
	\end{frame}
	
	\begin{frame}{Q35: Explanation}
		\begin{block}{Correct Answer: B}
			LLMs have demonstrated a willingness to share information they know is false and the ability to create new assumptions based on erroneous reasoning, referred to as "hallucinations."
		\end{block}
		\begin{block}{Option Analysis}
			\begin{itemize}
				\item \textbf{A - Incorrect:} LLMs can be inaccurate
				\item \textbf{B - Correct:} Hallucinations are a key risk
				\item \textbf{C - Incorrect:} Human oversight remains essential
				\item \textbf{D - Incorrect:} LLMs can be used maliciously
			\end{itemize}
		\end{block}
	\end{frame}
	
	\begin{frame}{Q36: AI Nontechnical Risk}
		\begin{block}{Question}
			Which is an example of a nontechnical risk that enterprises may overlook when using AI?
		\end{block}
		\begin{enumerate}[(A)]
			\item Software bugs
			\item Biased or incorrect outputs impacting human rights and safety
			\item Unpatched vulnerabilities
			\item Network latency
		\end{enumerate}
	\end{frame}
	
	\begin{frame}{Q36: Explanation}
		\begin{block}{Correct Answer: B}
			AI presents nontechnical risk, notably ethical and environmental impacts. Biased outputs that impact human rights and safety can result in reputational and legal damage.
		\end{block}
		\begin{block}{Option Analysis}
			\begin{itemize}
				\item \textbf{A - Incorrect:} Software bugs are technical risk
				\item \textbf{B - Correct:} Ethical/bias impacts
				\item \textbf{C - Incorrect:} Unpatched vulnerabilities are technical
				\item \textbf{D - Incorrect:} Network latency is technical
			\end{itemize}
		\end{block}
	\end{frame}
	
	\begin{frame}{Q37: Quantum Computing Impact}
		\begin{block}{Question}
			The MOST critical concern related to quantum computing for all organizations is its impact on:
		\end{block}
		\begin{enumerate}[(A)]
			\item Job displacement
			\item Applied cryptography
			\item Supply chain vulnerabilities
			\item Power generation sustainability
		\end{enumerate}
	\end{frame}
	
	\begin{frame}{Q37: Explanation}
		\begin{block}{Correct Answer: B}
			The most critical concern is quantum computing's impact on applied cryptography. Encryption used to secure transactions and sensitive information will effectively become obsolete once quantum computers become mainstream.
		\end{block}
		\begin{block}{Option Analysis}
			\begin{itemize}
				\item \textbf{A - Incorrect:} A concern but not the most universal
				\item \textbf{B - Correct:} Cryptography impact affects all
				\item \textbf{C - Incorrect:} A concern but not most critical
				\item \textbf{D - Incorrect:} A concern but not most critical
			\end{itemize}
		\end{block}
	\end{frame}
	
	\begin{frame}{Q38: Deepfake Social Engineering}
		\begin{block}{Question}
			Which emerging technology poses a substantial challenge for organizations that authorize financial transfers based on verbal approval?
		\end{block}
		\begin{enumerate}[(A)]
			\item Big data analytics
			\item Deepfake technology
			\item Blockchain
			\item Internet of Things
		\end{enumerate}
	\end{frame}
	
	\begin{frame}{Q38: Explanation}
		\begin{block}{Correct Answer: B}
			Deepfake technology can create convincing audio and video imitations, challenging organizations that authorize financial transfers based on verbal or visual approval. Real-world examples include wire transfers after deepfake video calls.
		\end{block}
		\begin{block}{Option Analysis}
			\begin{itemize}
				\item \textbf{A - Incorrect:} Not a deception technology
				\item \textbf{B - Correct:} Deepfakes enable impersonation
				\item \textbf{C - Incorrect:} Blockchain is a ledger technology
				\item \textbf{D - Incorrect:} IoT is about connected devices
			\end{itemize}
		\end{block}
	\end{frame}
	
	\begin{frame}{Q39: Blockchain Characteristic}
		\begin{block}{Question}
			Which characteristic makes it very difficult to falsify an entry in a blockchain?
		\end{block}
		\begin{enumerate}[(A)]
			\item Each entry is stored on a single central server
			\item Each entry contains both its own digital hash and the hash of the block before it
			\item Entries are encrypted with symmetric keys only
			\item Entries are never verified
		\end{enumerate}
	\end{frame}
	
	\begin{frame}{Q39: Explanation}
		\begin{block}{Correct Answer: B}
			The chain structure makes falsification difficult because each entry contains both its own hash and the hash of the previous block. Decentralization provides transparency.
		\end{block}
		\begin{block}{Option Analysis}
			\begin{itemize}
				\item \textbf{A - Incorrect:} Blockchain is decentralized
				\item \textbf{B - Correct:} Chained hashes
				\item \textbf{C - Incorrect:} Uses cryptographic hashing
				\item \textbf{D - Incorrect:} Entries are verified through consensus
			\end{itemize}
		\end{block}
	\end{frame}
	
	\begin{frame}{Q40: IoT Device Risk}
		\begin{block}{Question}
			Why do IoT devices pose a significant risk in enterprise environments?
		\end{block}
		\begin{enumerate}[(A)]
			\item They always have strong security features
			\item They tend to prioritize functionality over security, and their presence may become known only after deployment
			\item They are always managed by IT from the start
			\item They cannot be hacked
		\end{enumerate}
	\end{frame}
	
	\begin{frame}{Q40: Explanation}
		\begin{block}{Correct Answer: B}
			IoT devices tend to prioritize functionality over security. Their presence may become known to security staff only after deployment, leaving enterprises reliant on compensating controls.
		\end{block}
		\begin{block}{Option Analysis}
			\begin{itemize}
				\item \textbf{A - Incorrect:} IoT often lacks strong security
				\item \textbf{B - Correct:} Functionality over security, unknown deployment
				\item \textbf{C - Incorrect:} IoT often managed by facilities
				\item \textbf{D - Incorrect:} IoT devices can be hacked
			\end{itemize}
		\end{block}
	\end{frame}
	
	\begin{frame}{Q41: Omnipresent Connectivity Risk}
		\begin{block}{Question}
			Why does omnipresent connectivity make it harder to safeguard information?
		\end{block}
		\begin{enumerate}[(A)]
			\item People are less able to share information
			\item People are more able and more inclined to share information, including on social media
			\item Connectivity has no effect on information sharing
			\item All data is automatically encrypted
		\end{enumerate}
	\end{frame}
	
	\begin{frame}{Q41: Explanation}
		\begin{block}{Correct Answer: B}
			People are both more able to share information and more inclined to do so, including livestream video feeds taken without notice to those around them.
		\end{block}
		\begin{block}{Option Analysis}
			\begin{itemize}
				\item \textbf{A - Incorrect:} People are more able to share
				\item \textbf{B - Correct:} Increased sharing tendency
				\item \textbf{C - Incorrect:} Connectivity affects sharing
				\item \textbf{D - Incorrect:} Not all data is encrypted
			\end{itemize}
		\end{block}
	\end{frame}
	
	\begin{frame}{Q42: Big Data Challenge - Data Ownership}
		\begin{block}{Question}
			Which is a common challenge enterprises face associated with big data?
		\end{block}
		\begin{enumerate}[(A)]
			\item Data ownership is often not clearly defined
			\item Data is always perfectly organized
			\item Regulatory requirements are decreasing
			\item There is an abundance of skilled data analysts
		\end{enumerate}
	\end{frame}
	
	\begin{frame}{Q42: Explanation}
		\begin{block}{Correct Answer: A}
			A common challenge is that data ownership is often not clearly defined due to the misconception that data management is technical work. Other challenges include siloed structures, data deluge, lack of skilled analysts, and AI-related data biases.
		\end{block}
		\begin{block}{Option Analysis}
			\begin{itemize}
				\item \textbf{A - Correct:} Data ownership is often unclear
				\item \textbf{B - Incorrect:} Big data is often disorganized
				\item \textbf{C - Incorrect:} Data regulations are trending upward
				\item \textbf{D - Incorrect:} There is a shortage of skilled analysts
			\end{itemize}
		\end{block}
	\end{frame}
	
	% ================ FINAL SLIDES ================
	
	\begin{frame}{Technology Risk Key Takeaways}
		\begin{block}{Essential Concepts}
			\begin{itemize}
				\item EA provides current state, future vision, and a strategy to move between them
				\item Technology roadmaps link business goals to technology solutions
				\item Hardware risk: obsolescence, misconfiguration, default credentials
				\item OS risk: kernel-level vulnerabilities, unpatched flaws, end of support
				\item Cloud risk: reduced visibility, vendor lock-in, provider insiders
				\item IDS detects and alerts; IPS detects and blocks in-line
				\item Zero trust protects resources, not network segments
				\item Emerging tech introduces new risk categories
				\item Big data amplifies impact through centralization
			\end{itemize}
		\end{block}
	\end{frame}
	
	\begin{frame}{Final Advice}
		\begin{block}{Exam Preparation}
			\centering
			\vspace{0.3cm}
			\textbf{Remember:}
			\begin{itemize}
				\item Understand the \textit{why} behind each technology risk concept
				\item Know the difference between similar terms
				\item Understand how emerging technologies introduce new risk
				\item Practice scenario-based questions
				\item Use the ISACA mindset
			\end{itemize}
			\vspace{0.5cm}
			\textbf{Good Luck on Your CRISC Exam!}
		\end{block}
	\end{frame}
	
\end{document}